\section{Regularity of the finite root-stack presentations}
\label{app:regular-presentation}

This appendix proves the two statements on regularity summarized at the end
of Section~\ref{subsec:root-stage-regularity}, whose notation we keep.  Thus
$\mathbb T=\Spec\Bbbk[\widetilde M]$ is the character torus in the frame of the
sheet $S_0$, and $T\subset\Tcal$ is a finite rooted subtree.  At a Kummer edge
$e$ of $T$ we have the root degree $d_e\in\{2,4\}$, the transported radicand
$\varrho_e$, the Laurent polynomial $\Phi_e$ of
Lemma~\ref{lem:transported-radicands} and the divisor
$D_e=\operatorname{div}\Phi_e$ of zeros of $\varrho_e$; moreover
$D_T=\bigcup_eD_e$, the scheme $P_T$ is the Kummer cover, with
group $G_T=\prod_e\mu_{d_e}$, so that $\Xfrak_T=[P_T/G_T]$, and $P_T^\circ$
is the preimage of $\mathbb T\setminus D_T$ in $P_T$.

\subsection{A multiple-root obstruction}

\begin{proposition}[A multiple-root obstruction]
\label{prop:root-stage-nonregular}
Write $f_B(x)=1+a_\rho x+x^2=(1+\lambda x)(1+\lambda^{-1}x)$ for the $B_0$-slab
function, with $\lambda\neq\pm1$, and let
\[
 \varrho_n=f_B\bigl(u^nx_B\bigr),\qquad u=\bigl(1+x_A^{-1}\bigr)^4,\qquad n\in\ZZ,
\]
be the transported $B_0$-radicands of Theorem~\ref{thm:infinite-root-orbit},
adjoined at the crossings $e_n$ of the $B_0$-arc that follow the paths
$\gamma_{\widetilde A}^n$.  Let $T\subset\Tcal$ contain
$e_0,e_1,e_2$, and let $p\in\mathbb T$ be a point with
$(x_A,x_B)(p)=(-\tfrac12,-\lambda^{-1})$.  Then the Kummer cover $P_T$ is not
regular at any point $\hat p$ over $p$.  At $\hat p$ it is not log regular,
either with the trivial log structure or with the log structure given by
the chart from the free commutative monoid on the Kummer edges of $T$ that
sends each Kummer edge $e$ to $g_e$.  In particular $\Xfrak_T$ is not regular.
\end{proposition}

\begin{proof}
The characters $x_A=z^{2u_{\widetilde A}}$ and $x_B=z^{2u_{B_0}}$ have
linearly independent exponents, so $(x_A,x_B):\mathbb T\to\mathbb G_m^2$ is
\'etale, and $x_A+\frac12$, $x_B+\lambda^{-1}$ form a regular system of
parameters of $\Ocal_{\mathbb T,p}$; we compute differentials in these coordinates.
Put $\varrho_n^\pm=1+\lambda^{\pm1}u^nx_B$, so that
$\varrho_n=\varrho_n^+\varrho_n^-$.  Then $u(p)=1$ and $f_A(p)=\tfrac12\neq0$,
and for every $n$
\[
 \varrho_n^+(p)=0,\qquad \varrho_n^-(p)=1-\lambda^{-2}\neq0,\qquad
 \nabla\varrho_n^+(p)=(-16n,\lambda).
\]
Near $p$ each $\varrho_n$ is a unit times $\varrho_n^+$.  For $n=0,1,2$ the
gradients $(0,\lambda),(-16,\lambda),(-32,\lambda)$ are pairwise
independent, so the three root divisors are smooth and pairwise transverse
at $p$.

For $n\geq0$ the radicand $\varrho_n$ is a Laurent polynomial, so we may
take $\Phi_{e_n}=\varrho_n$ in Lemma~\ref{lem:transported-radicands}; the
fixed $\Phi_{e_n}$ differs from $\varrho_n$ by the fourth power of a unit
$v$, which changes $P_T$ by the isomorphism $g\mapsto vg$.  Write
$g_0,g_1,g_2$ for the roots at $e_0,e_1,e_2$, of degree four, and
$g_3,\ldots,g_k$ for the roots at the other Kummer edges of $T$, with $d_j$
and $\Phi_j$ accordingly.  Then
$P_T=\Spec\Ocal_{\mathbb T}[g_0,\ldots,g_k]/(g_j^{d_j}-\Phi_j)$ is finite and flat
over the surface $\mathbb T$, hence of dimension two at every point, and
$g_0,g_1,g_2$ vanish at every point $\hat p$ over $p$.  The Jacobian of the
defining equations at $\hat p$ has rank at most $2+(k-2)=k$, because the
rows of $g_j^4-\varrho_j$, $j\leq2$, vanish in the $g$-columns.  The tangent
space at $\hat p$ therefore has dimension at least $(k+3)-k=3>2$, and $P_T$
is not regular at $\hat p$.  With the trivial log structure, log regularity
at $\hat p$ is regularity \cite[Def.~2.1]{kato1994toric}, so it fails.  With
the log structure of the roots, the characteristic monoid at $\hat p$ is
free on the Kummer edges whose roots vanish at $\hat p$, so its rank is at
least three, and the ideal $\mathfrak a_{\hat p}\subset\Ocal_{\hat p}$
generated by the images of the nonunits of the log structure contains
$g_0,g_1,g_2$.  The quotient $\Ocal_{\hat p}/\mathfrak a_{\hat p}$ is finite
over $\Ocal_{\mathbb T,p}/(\varrho_0^+,\varrho_1^+)=\Bbbk$, because $\varrho_0^+$ and
$\varrho_1^+$ vanish transversely at $p$, and so has dimension zero.  Kato's
condition
$\dim\Ocal_{\hat p}=\dim(\Ocal_{\hat p}/\mathfrak a_{\hat p})+\operatorname{rank}$
\cite[Def.~2.1]{kato1994toric} would read $2=0+\operatorname{rank}$ with a
rank of at least three, so it fails.  Since $P_T\to\Xfrak_T$ is \'etale and
surjective, $\Xfrak_T$ is not regular.  A tree of this kind exists: by
Theorem~\ref{thm:infinite-root-orbit} the radicands
$\varrho_0,\varrho_1,\varrho_2$ occur along the finitely many paths
$\gamma_{\widetilde A}^n$, $n=0,1,2$, followed by a crossing of the
$B_0$-arc, and their rooted union is a finite subtree.
\end{proof}

At the cusps $a_\rho=\pm2$ the same point, with $\lambda=\pm1$, gives the same
conclusion.  There $\varrho_n=(\varrho_n^+)^2$ by \eqref{eq:cusp-roots}, and
both the fourth roots of $\varrho_n$ and the square roots of $\varrho_n^+$,
whose gradients at $p$ are again $(-16n,\lambda)$, give Kummer covers whose
tangent space at the points over $p$ has dimension at least three.

\subsection{A regular dominating presentation}

Regularity is recovered by a proper modification of the finite
presentations, together with the modification of their Kummer covers that
it induces; the latter carries the lifted transitions over the preimage of
$\mathbb T\setminus D_T$ (Section~\ref{subsec:root-stage-regularity}).  A
reduced curve on a smooth surface is a \emph{simple normal crossings
divisor} if its components are smooth, meet pairwise
transversely, and no three pass through a point.  The modification is the
stack of fourth roots of the boundary components of the minimal log
resolution of $(\mathbb T,D_T)$; fourth roots suffice because every root
degree $d_e$ divides $4$.

\begin{theorem}[Regular dominating presentation]
\label{thm:regular-presentation}
For every finite rooted subtree $T\subset\Tcal$ there are a smooth
separated Deligne--Mumford stack $\mathcal Y_T$ over $\Bbbk$, a divisor
$\mathcal E_T\subset\mathcal Y_T$ that is a simple normal crossings divisor
on an \'etale atlas, and a proper morphism
$p_T:\mathcal Y_T\to\Xfrak_T$.  Put
$\widetilde{\mathcal Y}_T=\mathcal Y_T\times_{\Xfrak_T}P_T$, let
$\widetilde p_T:\widetilde{\mathcal Y}_T\to P_T$ be the projection, and let
$\widetilde{\mathcal E}_T$ be the preimage of $\mathcal E_T$.  Then the
following hold.
\begin{enumerate}
\item $p_T$ is an isomorphism over $\mathbb T\setminus D_T$.  The projection
$\widetilde{\mathcal Y}_T\to\mathcal Y_T$ is a $G_T$-torsor, so
$\widetilde{\mathcal Y}_T$ is a smooth separated Deligne--Mumford stack and
$\mathcal Y_T=[\widetilde{\mathcal Y}_T/G_T]$.  The morphism $\widetilde p_T$
is proper and $G_T$-equivariant, and it is an isomorphism over
$P_T^\circ$.  Consequently every function, section of $\Lcal_\kappa$,
bracket or automorphism defined over
$P_T^\circ\times_{\Spec\Bbbk}\Spec R_{\Wcal,N}$ pulls back unchanged to the
preimage of $\mathbb T\setminus D_T$ in
$\widetilde{\mathcal Y}_T\times_{\Spec\Bbbk}\Spec R_{\Wcal,N}$; this applies
in particular to the lifted transitions along $T$, to the Jacobi bracket of
Proposition~\ref{prop:Jacobi-transport}, and to the wall automorphisms
base-changed along $\upsilon_\Wcal$ of \eqref{eq:window-map}.
\item For every labelled truncation $\Wcal$ and every $N\geq1$, the base
changes of $\mathcal Y_T$ and of $\widetilde{\mathcal Y}_T$ to
$\Spec R_{\Wcal,N}$, with
the log structures of $\mathcal E_T$ and of $\widetilde{\mathcal E}_T$, are
flat and log smooth over $\Spec R_{\Wcal,N}$ with the trivial log
structure, and they have smooth atlases over $R_{\Wcal,N}$.  Their closed
fibers $(\mathcal Y_T,\mathcal E_T)$ and
$(\widetilde{\mathcal Y}_T,\widetilde{\mathcal E}_T)$ are log regular.
\item If $T\subset T'$, there is a proper morphism
$q_{T'T}:\mathcal Y_{T'}\to\mathcal Y_T$ with
$p_T\circ q_{T'T}\simeq\pi_{T'T}\circ p_{T'}$, where
$\pi_{T'T}:\Xfrak_{T'}\to\Xfrak_T$ forgets the roots along the divisors
indexed by $T'\setminus T$.  It lifts to a proper morphism
$\widetilde q_{T'T}:\widetilde{\mathcal Y}_{T'}\to\widetilde{\mathcal Y}_T$
over the projection $P_{T'}\to P_T$ that forgets the same roots.  All
constructions commute with the base changes $R_{\Wcal',N'}\to R_{\Wcal,N}$.
\end{enumerate}
\end{theorem}

\begin{proof}
\emph{Embedded resolution.}  Let $\pi_T:\widehat{\mathbb T}_T\to\mathbb T$ be the minimal log
resolution of $(\mathbb T,D_T)$, obtained by repeatedly blowing up the
finitely many points at which the reduced total transform of $D_T$ is
not a simple normal crossings divisor, and let
$E_T=E_1\cup\cdots\cup E_{n_T}$ be the reduced total transform of
$D_T$ in $\widehat{\mathbb T}_T$.  The process stops, and every composition of point
blowups with centers over $\mathbb T$ under which the reduced total transform of
$D_T$ is a simple normal crossings divisor factors through $\pi_T$.
Indeed, on a smooth projective compactification of $\mathbb T$ such compositions
exist by embedded resolution of curves on surfaces
\cite[Thm.~V.3.9]{hartshorne1977}, followed if necessary by the blowups of
the points where three components meet.  If such a composition factors
through a stage of the process, it is not an isomorphism over any
neighborhood of a center of the next step, since the reduced total
transform of $D_T$ is not a simple normal crossings divisor there, so
its inverse is not defined at that center; by
\cite[Prop.~V.5.3]{hartshorne1977}, applied after blowing up the same centers
in a smooth projective compactification of $\mathbb T$, it factors through the blowup
of each such center, hence through the next stage.  As each stage adds an
exceptional curve of the composition, the process stops.  All centers lie
over $D_T$, so $\pi_T$ is an isomorphism over $\mathbb T\setminus D_T$.

\emph{Fourth roots of the boundary.}  Let $\mathcal Y_T$ be the fiber
product over $\widehat{\mathbb T}_T$ of the root stacks $\sqrt[4]{(\widehat{\mathbb T}_T,E_j)}$, and let
$\mathcal E_j$ be the reduced preimage of $E_j$, so that $E_j$ pulls back
to $4\mathcal E_j$.  Where $E_{j_1},\ldots,E_{j_l}$ meet, with local
equations $x_1,\ldots,x_l$ forming part of a regular system of parameters,
$\mathcal Y_T$ has the \'etale chart
$U=\Spec\Ocal[y_1,\ldots,y_l]/(y_1^4-x_1,\ldots,y_l^4-x_l)$ with its action
of $\mu_4^l$ \cite[Ex.~3.1]{cadman2007}; it is smooth over $\Bbbk$, and
$\mathcal E_T=\bigcup_j\mathcal E_j$ is $y_1\cdots y_l=0$ on it.  Since
$\mu_4$ is \'etale in characteristic zero, $\mathcal Y_T$ is a smooth
separated Deligne--Mumford stack and $\mathcal E_T$ is a simple normal
crossings divisor on an \'etale atlas.

\emph{The morphism $p_T$.}  Write $\pi_T^*D_e=\sum_ja_{ej}E_j$
with $a_{ej}\geq0$.  On $\mathcal Y_T$ this divisor is $d_e$
times $\sum_j(4a_{ej}/d_e)\mathcal E_j$, which is integral
because $d_e$ divides $4$; the line bundle of the latter, its
canonical section and the canonical isomorphism of its $d_e$th power
with $\Ocal(\pi_T^*D_e)$ define a morphism
$\mathcal Y_T\to\sqrt[d_e]{(\mathbb T,D_e)}$ over $\pi_T$
\cite[Defs.~2.1, 2.3]{cadman2007}.  Together these give $p_T$, which is
proper because $\mathcal Y_T\to\mathbb T$ is proper and $\Xfrak_T\to\mathbb T$ is
separated.  Locally $\pi_T^*\Phi_e=v_e\prod_jx_j^{a_{ej}}$
with $v_e$ a unit, so after the \'etale extension
$\eta_e^{d_e}=v_e$ a $d_e$th root of
$\pi_T^*\Phi_e$ is a unit times a monomial in the boundary
coordinates:
\begin{equation}\label{eq:resolved-root-formula}
 g_e=\eta_e
 \prod_{j}y_j^{4a_{ej}/d_e},
 \qquad
 g_e^{d_e}=\pi_T^*\Phi_e .
\end{equation}

\emph{The Kummer cover.}  As a base change of the $G_T$-torsor
$P_T\to\Xfrak_T$, with $G_T$ finite and \'etale,
$\widetilde{\mathcal Y}_T\to\mathcal Y_T$ is a finite \'etale $G_T$-torsor.
Hence $\mathcal Y_T=[\widetilde{\mathcal Y}_T/G_T]$, the stack
$\widetilde{\mathcal Y}_T$ is smooth, separated and Deligne--Mumford, and
$\widetilde p_T$, the base change of $p_T$, is proper and
$G_T$-equivariant.  Over a chart $U$ as above,
$\widetilde U=U\times_{\mathcal Y_T}\widetilde{\mathcal Y}_T$ is a scheme,
finite and \'etale over $U$ and hence smooth over $\Bbbk$, and
$\widetilde{\mathcal E}_T$ is $y_1\cdots y_l=0$ on it.  On $\widetilde U$
the quotient $g_e\prod_jy_j^{-4a_{ej}/d_e}$ of the root
$g_e$ by the monomial is a rational function whose $d_e$th power
is the unit $v_e$, hence a unit because $\widetilde U$ is normal, and
\eqref{eq:resolved-root-formula} holds on $\widetilde U$.

\emph{Proof of~(1) and~(2).}  No divisor that is blown up or rooted meets
$\mathbb T\setminus D_T$, so $p_T$ is an isomorphism over $\mathbb T\setminus D_T$
and $\widetilde p_T$ is one over $P_T^\circ$; this proves~(1).  The chart
$\NN^l\to\Ocal_U$ that sends the $j$th basis vector to $y_j$ induces smooth
morphisms from $U$ and from $\widetilde U$, which is \'etale over $U$, to
$\Spec\Bbbk[\NN^l]$, so both are log smooth over $\Bbbk$ by Kato's criterion
\cite[Thm.~3.5]{kato1989log}; log smoothness, flatness and smooth atlases
persist after base change to $R_{\Wcal,N}$.  At a point on exactly $l$
components, the characteristic monoid is $\NN^l$, the ideal generated by
its nonunits is $(y_1,\ldots,y_l)$, and the quotient by this ideal is
regular of dimension $2-l$; this is Kato's condition for log regularity
\cite[Def.~2.1]{kato1994toric}, which proves~(2).

\emph{Proof of~(3).}  Since $D_T\subset D_{T'}$, the reduced
total transform of $D_T$ in $\widehat{\mathbb T}_{T'}$ is a union of components of
$E_{T'}$, hence a simple normal crossings divisor.  By minimality
$\pi_{T'}=\pi_T\circ b_{T'T}$ for a morphism $b_{T'T}:\widehat{\mathbb T}_{T'}\to\widehat{\mathbb T}_T$, and
since $b_{T'T}^{-1}(E_T)\subset E_{T'}$, each $E_j$ pulls back to
$\sum_kc_{jk}E'_k$ with $c_{jk}\geq0$, which is
$4\sum_kc_{jk}\mathcal E'_k$ on $\mathcal Y_{T'}$.  As for $p_T$, this
defines $q_{T'T}$ over $b_{T'T}$, proper because $\mathcal Y_{T'}\to\widehat{\mathbb T}_T$
is proper and $\mathcal Y_T\to\widehat{\mathbb T}_T$ is separated.  For each Kummer edge $e$
of $T$, both $p_T\circ q_{T'T}$ and $\pi_{T'T}\circ p_{T'}$ correspond, on
the factor $\sqrt[d_e]{(\mathbb T,D_e)}$, to the line bundle
$\Ocal(\sum_k(4a'_{ek}/d_e)\mathcal E'_k)$ with its canonical
section and isomorphism, where $a'_{ek}$ is the multiplicity of
$\pi_{T'}^*D_e$ along $E'_k$; hence they are $2$-isomorphic.  The
projection $P_{T'}\to P_T$ is equivariant for $G_{T'}\to G_T$ and lies over
$\pi_{T'T}$; with $q_{T'T}$ and this $2$-isomorphism it induces
$\widetilde q_{T'T}$, which is proper because
$\widetilde{\mathcal Y}_{T'}\to\mathbb T$ is proper and
$\widetilde{\mathcal Y}_T\to\mathbb T$ is separated.  Every object is a base
change from $\Bbbk$, which gives the last assertion.
\end{proof}
